\documentclass[11pt]{beamer}

% ---------------------------------------------------------------------
% Preamble
% ---------------------------------------------------------------------

% ---------------------------------------------------------------------
% Packages
% ---------------------------------------------------------------------

\usepackage[utf8]{inputenc}
\usepackage{amsmath}
\usepackage{amssymb}
\usepackage[ngerman]{babel}
\usepackage{beamerprosper}
\usepackage{color}
\usepackage{epsfig}
\usepackage{graphicx}
\usepackage{epsfig}
\usepackage{latexsym}
\usepackage{listings}
\usepackage{times}
\usepackage{url}
\usepackage{xspace}
\usepackage{pgfpages}

% ---------------------------------------------------------------------
% Spacings
% ---------------------------------------------------------------------

\setlength{\parskip}{0.5ex}

% ---------------------------------------------------------------------
% Settings
% ---------------------------------------------------------------------

\lstset{ %
  backgroundcolor=\color{white},   % choose the background color; you must add \usepackage{color} or \usepackage{xcolor}
  basicstyle=\footnotesize\ttfamily, % the size of the fonts that are used for the code
  breakatwhitespace=false,         % sets if automatic breaks should only happen at whitespace
  breaklines=true,                 % sets automatic line breaking
  captionpos=b,                    % sets the caption-position to bottom
  commentstyle=\color{green},  	   % comment style
  deletekeywords={...},            % if you want to delete keywords from the given language
  escapeinside={\%*}{*)},          % if you want to add LaTeX within your code
  extendedchars=true,              % lets you use non-ASCII characters; for 8-bits encodings only, does not work with UTF-8
  frame=single,                    % adds a frame around the code
  keepspaces=true,                 % keeps spaces in text, useful for keeping indentation of code (possibly needs columns=flexible)
  keywordstyle=\color{blue},       % keyword style
  language=Java,                   % the language of the code
  morekeywords={*,...},            % if you want to add more keywords to the set
  numbers=left,                    % where to put the line-numbers; possible values are (none, left, right)
  numbersep=5pt,                   % how far the line-numbers are from the code
  numberstyle=\tiny\color{green},  % the style that is used for the line-numbers
  rulecolor=\color{black},         % if not set, the frame-color may be changed on line-breaks within not-black text (e.g. comments (green here))
  showspaces=false,                % show spaces everywhere adding particular underscores; it overrides 'showstringspaces'
  showstringspaces=false,          % underline spaces within strings only
  showtabs=true,                   % show tabs within strings adding particular underscores
  stepnumber=1,                    % the step between two line-numbers. If it's 1, each line will be numbered
  stringstyle=\color{blue},        % string literal style
  tabsize=2,                       % sets default tabsize to 2 spaces
  title=\lstname                   % show the filename of files included with \lstinputlisting; also try caption instead of title
}

% ---------------------------------------------------------------------
% Paths
% ---------------------------------------------------------------------

\graphicspath{{./images/}}

% ---------------------------------------------------------------------
% Theme
% ---------------------------------------------------------------------

\mode<all>
\usetheme{CambridgeUS}
\mode<presentation>
{
  \useinnertheme{rectangles}
  %\useoutertheme{}
  \setbeamercovered{transparent}
  \setbeamertemplate{navigation symbols}{}
}

% \definecolor{darkred}{rgb}{0.625,0.125,0.25}
\definecolor{darkred}{rgb}{0,0.4196,0.5803}
% \definecolor{bauhausblue}{rgb}{0,0.107,0.148}

\setbeamercolor{block body}{bg=gray!20}
\setbeamercolor{block title}{bg=darkred,fg=white}
\setbeamercolor{footlinecolorl}{fg=black,bg=lightgray}
\setbeamercolor{footlinecolor}{fg=black,bg=gray}
\setbeamercolor{footlinecolord}{fg=black,bg=darkgray}

\setbeameroption{show notes}
\setbeameroption{show notes on second screen=right}

% Auskommentieren, für Section 1,2,3, ...
\setbeamertemplate{section page}
{
    \begin{centering}
    \begin{beamercolorbox}[sep=12pt,center]{part title}
    \usebeamerfont{section title}\insertsection\par
    \end{beamercolorbox}
    \end{centering}
}

\setbeamertemplate{footline}{%
	\hbox{%
	\begin{beamercolorbox}[wd=.40\paperwidth,ht=4.25ex,left,leftskip=3ex]{author in head/foot}%
	    \vbox to4.25ex{\vfil\hbox{\usebeamerfont{author in head/foot} \insertshortauthor}\vfil}%
	\end{beamercolorbox}%
	\begin{beamercolorbox}[wd=.30\paperwidth,ht=4.25ex,center]{title in head/foot}%
	    \vbox to4.25ex{\vfil\hbox{\usebeamerfont{date in head/foot}\insertshorttitle{}}\vfil}%
	\end{beamercolorbox}%
  % \begin{frame number}{}
  % \end{frame number}%
	\begin{beamercolorbox}[wd=.30\paperwidth,ht=4.25ex,right,rightskip=3ex]{date in head/foot}%
	    % \vbox to4.25ex{\vfil\hbox{\insertshortdate{}}\vfil}%
      % \vbox to4.25ex{\vfill\hbox{frame number}}
      \vbox to4.25ex{\vfill\hbox{\insertframenumber{} / \inserttotalframenumber }\vfill}
	\end{beamercolorbox}}%
}
%\setbeamertemplate{footline}{
%  \quad \tiny \insertshortauthor \hfill \insertshorttitle \qquad \hfill \insertshortdate\ \qquad\qquad
%}
\setbeamertemplate{theorems}[numbered]
% ---------------------------------------------------------------------
% Commands
% ---------------------------------------------------------------------

\newcommand{\eg}{\textit{e.g.,}\xspace}
\newcommand{\shortsep}{||}
\newcommand{\sep}{\ \shortsep \ }
\def\getpdfpages#1#2{\begingroup
  \setbeamercolor{background canvas}{bg=}
  \includepdf[pages={#1},%
  pagecommand={\global\setcounter{framenumber}{\value{page}}%
    \expandafter\def\expandafter\insertshorttitle\expandafter{%
      \insertshorttitle\hfill\insertframenumber\,/\,\inserttotalframenumber}}%
  ]{#2}
  \endgroup}

% ---------------------------------------------------------------------
% Title
% ---------------------------------------------------------------------

\title[Bauhausboards]{Bauhausboards}
\subtitle{Interactive Door Signs for the Office}
\author[André Karge]{André Karge}
\institute[Bauhaus-Universität Weimar]{}
\date[\today]{\today}

\raggedright
\AtBeginSection{\frame{\sectionpage}}

% ---------------------------------------------------------------------

\begin{document}

% ---------------------------------------------------------------------
% Content
% ---------------------------------------------------------------------

\maketitle

\begin{frame}{Agenda}
\tableofcontents
\end{frame}

\section{Einleitung}
\section{Verwandte Arbeiten}
\section{Umsetzung}
\section{Studie}
\section{Ausblick}


% % ---------------------------------------------------------------------
% \section{Aufbau eines Vortrags}
% % ---------------------------------------------------------------------

% \begin{frame}{Aufbau eines Vortrags}
%   \begin{itemize}
%     \item Motivation des Themas
%     % \pause
%     \item Erläuterung der Problemstellung
%     % \pause
%     \item Bisherige/Mögliche Ansätze
%     \begin{itemize}
%         \item Ansatz 1
%         \item Ansatz 2
%     \end{itemize}
%     \item Eigene Lösung oder -versuche
%     \item Zusammenfassung oder Schlussfolgerung
%   \end{itemize}
% \end{frame}

% % ---------------------------------------------------------------------
% \section{Layouts}
% % ---------------------------------------------------------------------

% \begin{frame}{Block- und Mathe-Umgebungen}
% \begin{block}{Definition}
% Dies ist eine Definition.
% $$ a^2 + b^2 = c^2 $$
% \end{block}
% Der erste Buchstabe im griechischen Alphabet ist $\alpha$.
% Grafiken lassen sich einfach mit \texttt{\textbackslash includegraphics} einfügen.
% \includegraphics[width=0.7\textwidth]{buw-logo}
% \end{frame}

% % ---------------------------------------------------------------------

% \begin{frame}{Mehrspaltig}
%   Mit Hilfe von \texttt{\textbackslash columns} lassen sich mehrspaltige Folien erstellen.
%   \begin{columns}
%     \begin{column}{0.32\textwidth}
%         \includegraphics[width=0.9\textwidth]{Tux}
%     \end{column}
%     \begin{column}{0.32\textwidth}
%         \includegraphics[width=0.9\textwidth]{Tux_ecb}
%     \end{column}
%     \begin{column}{0.32\textwidth}
%         \includegraphics[width=0.9\textwidth]{Tux_secure}
%     \end{column}
%   \end{columns}
% \end{frame}

% % ---------------------------------------------------------------------

% \begin{frame}{Tabellen}
%   Tabellen sind auch ganz einfach.
%   \begin{table}
%     \begin{tabular}{lcr}
%         \hline
%         Linksbündig & Zentriert & Rechtsbündig \\
%         \hline
%         1 & 2 & 3 \\
%         4 & 5 & 6 \\
%         \hline
%     \end{tabular}
%     \caption{Tabellenunterschrift}
%   \end{table}
% \end{frame}

% % ---------------------------------------------------------------------

% \begin{frame}{Weitere Informationen}
%   Homepage: \\
%   \url{https://bitbucket.org/rivanvx/beamer/wiki/Home}
  
%   Benutzerhandbuch: \\
%   \url{http://www.ctan.org/tex-archive/macros/latex/contrib/beamer/doc/beameruserguide.pdf}
  
%   Wiki LaTeX-Kompendium: \\
%   \url{http://de.wikibooks.org/wiki/LaTeX-Kompendium}
  
%   Stack Exchange für LaTeX: \\
%   \url{http://tex.stackexchange.com/}
  
%   Symbolsuche: \\
%   \url{http://detexify.kirelabs.org/classify.html}
% \end{frame}

% % ---------------------------------------------------------------------

% \begin{frame}[fragile]{Listing}
% \begin{lstlisting}
% public class Calculator {

%     private int memory;

%     public int add(int a, int b) {
%         String foo = "bar"; // Just to demonstrate strings
%         return a + b;
%     }

%     /**
%      * Very important function!
%      * Good comments are essential to deliver high-quality code!
%      */
%     public int multiply(int a, int b) {
%         return a * b;
%     }
% }
% \end{lstlisting}
% \end{frame}

% % ---------------------------------------------------------------------
% % Appendix
% % ---------------------------------------------------------------------

% \appendix
\end{document}
